\chapter{Radical Prostatectomy}

\section*{Introduction \& Clinical Indications}
A radical prostatectomy is a definitive surgical procedure aimed at the complete excision of the prostate gland, seminal vesicles, and often the surrounding pelvic lymph nodes. This operation is primarily indicated for patients diagnosed with clinically localized prostate cancer, where the malignancy is confined to the prostate or has minimally extended beyond its capsule. The significance of this surgery lies in its potential to eradicate cancerous tissues while striving to preserve critical anatomical structures that are essential for urinary continence and erectile function. The procedure is tailored to individual patient characteristics, including cancer staging and overall health, and can be performed using various approaches, including open, laparoscopic, or robotic-assisted techniques. The ultimate goal is to achieve optimal oncologic outcomes while maintaining the patient's quality of life.

\begin{itemize}
  \item Open Radical Retropubic Prostatectomy (RRP)
  \item Laparoscopic Radical Prostatectomy (LRP)
  \item Robotic-Assisted Laparoscopic Prostatectomy (RALP)
  \item Prostate Cancer Surgery
  \item Nerve-Sparing Prostatectomy
  \item Total Prostatectomy
\end{itemize}

The surgical principle of radical prostatectomy is the complete en bloc excision of the prostate gland, seminal vesicles, and any involved lymph nodes. This approach aims to achieve maximal oncologic control by removing all known and potential microscopic disease within these structures. The rationale behind this procedure is to ensure negative surgical margins, which indicates that no cancer cells are present at the edges of the excised tissue, thereby minimizing the risk of local recurrence.

A critical aspect of the surgery is the meticulous dissection around the neurovascular bundles to preserve erectile function and the external urethral sphincter to maintain urinary continence. The technical goals of the procedure include definitive cancer eradication and minimizing functional morbidity related to urinary and sexual function. After the prostate is removed, the bladder neck is reconstructed, and the remaining urethra is anastomosed to the bladder neck, creating a new, watertight connection for urine outflow. In select cases, pelvic lymph node dissection is performed to accurately stage the disease and potentially remove micrometastases, further enhancing oncologic control, particularly in patients with higher-risk disease features.

The history of radical prostatectomy dates back to the late 19th and early 20th centuries, with early attempts often resulting in significant morbidity. In 1904, Hugh Hampton Young performed the first successful radical perineal prostatectomy at Johns Hopkins, marking a pivotal moment in the surgical treatment of prostate cancer. However, this technique was associated with high rates of urinary incontinence and was technically challenging, limiting its adoption.

A significant advancement occurred in 1947 when Terence Millin introduced the radical retropubic prostatectomy, which improved surgical exposure and facilitated pelvic lymph node dissection. This open approach became the standard for decades. The most transformative development in modern radical prostatectomy came in the early 1980s with Patrick Walsh's detailed anatomical studies, which identified the neurovascular bundles and external urethral sphincter. His work on the "anatomic radical prostatectomy" significantly reduced rates of incontinence and erectile dysfunction, making the procedure more acceptable to patients.

The late 1990s saw the introduction of laparoscopic radical prostatectomy, and by the early 2000s, robotic-assisted laparoscopic prostatectomy (RALP) emerged as the dominant surgical approach, offering advantages in visualization, dexterity, and potentially faster recovery times.

Radical prostatectomy is indicated for men diagnosed with clinically localized prostate cancer who meet specific criteria for curative intervention.

\begin{itemize}
  \item Clinically localized prostate cancer: Cancer confined to the prostate gland or with minimal extracapsular extension, as determined by biopsy and imaging.
  \item Favorable life expectancy: Patients with a life expectancy of at least 10-15 years to benefit from long-term cancer control.
  \item Intermediate- or high-risk prostate cancer: Patients with Gleason scores of 7 or higher, elevated PSA levels, or palpable tumors where active surveillance is not appropriate.
  \item Patient preference: After thorough counseling on treatment options, patients may choose surgery as their preferred definitive treatment.
  \item Younger age: Generally favored in younger men who can tolerate surgery and benefit from long-term cancer control.
  \item High-volume low-risk disease: Some patients with extensive but low-grade disease may opt for surgery to avoid uncertainties associated with active surveillance.
  \item Localized recurrence after prior local therapy: Considered as salvage therapy in selected cases after failed radiation therapy, though this carries higher risks.
\end{itemize}

Radical prostatectomy is primarily considered a first-line curative surgical treatment for localized prostate cancer. It stands as one of the main definitive treatment options alongside external beam radiation therapy and brachytherapy for men with intermediate- or high-risk localized disease and a good life expectancy. The primary role of radical prostatectomy is to remove the entire tumor burden with the intent of achieving a complete cure and long-term cancer control.

In certain scenarios, radical prostatectomy may serve as a secondary or salvage procedure, particularly in patients who have undergone prior local therapy, such as radiation, and have experienced local recurrence. Salvage prostatectomy is technically more challenging due to tissue fibrosis and altered anatomy, increasing the risk of complications. Therefore, its application is reserved for highly selected patients in specialized centers.

\section*{Relevant Surgical Anatomy}
The prostate gland is a walnut-sized exocrine gland located inferior to the bladder and anterior to the rectum, encircling the prostatic urethra. It is anatomically divided into several zones: the peripheral zone, which is the most common site for cancer; the central zone; and the transitional zone, which is often involved in benign prostatic hyperplasia. The prostate is connected superiorly to the bladder neck and inferiorly to the external urethral sphincter, which plays a crucial role in urinary continence.

Adjacent to the prostate are the seminal vesicles, which produce seminal fluid, and the vas deferens, which transports sperm. These structures are closely associated with the prostate and are typically excised during radical prostatectomy due to the risk of cancer spread. The neurovascular bundles, which are essential for erectile function, run posterolaterally to the prostate and are located within Denonvilliers' fascia. Preservation of these bundles is a key consideration during surgery. 

The arterial supply to the prostate is primarily from the inferior vesical artery, a branch of the internal iliac artery, while venous drainage occurs through the prostatic venous plexus into the internal iliac veins. Lymphatic drainage follows the internal iliac, obturator, and external iliac lymph nodes, which are often sampled during pelvic lymph node dissection.

\begin{itemize}
  \item Prostate gland: walnut-sized, located inferior to the bladder.
  \item Seminal vesicles: produce seminal fluid, closely associated with the prostate.
  \item Neurovascular bundles: critical for erectile function, located posterolaterally.
  \item Arterial supply: inferior vesical artery from the internal iliac artery.
  \item Venous drainage: prostatic venous plexus into internal iliac veins.
  \item Lymphatic drainage: follows internal iliac, obturator, and external iliac nodes.
\end{itemize}

\section*{Preoperative Workup \& Preparation}
Thorough preoperative preparation is crucial for a safe and successful radical prostatectomy. This comprehensive assessment ensures the patient is an appropriate candidate and is adequately prepared for surgery.

\begin{itemize}
  \item **Diagnostic Workup:** Confirmation of prostate cancer diagnosis via biopsy, including Gleason score and tumor volume. Preoperative staging with PSA levels, digital rectal exam, and imaging (e.g., multiparametric MRI) to assess for local extension or metastasis.
  
  \item **Medical Evaluation:** Comprehensive assessment of the patient's overall health, including cardiac, pulmonary, and renal function, involving blood tests and an electrocardiogram.
  
  \item **Medication Review:** All medications, especially blood thinners, must be reviewed and typically discontinued several days prior to surgery to minimize bleeding risk.
  
  \item **Bowel Preparation:** A bowel preparation regimen may be prescribed to reduce the risk of rectal injury and facilitate surgical exposure.
  
  \item **NPO Status:** Patients are instructed to be NPO for a specified period before surgery to prevent aspiration during anesthesia.
  
  \item **Preoperative Counseling:** Detailed discussion with the surgeon regarding the procedure, potential risks, expected outcomes, and alternative treatments. Informed consent is obtained.
  
  \item **Antibiotic Prophylaxis:** Intravenous antibiotics are administered shortly before the incision to reduce the risk of surgical site infection.
  
  \item **Pelvic Floor Exercises:** Patients may be advised to begin pelvic floor exercises preoperatively to strengthen these muscles, aiding in recovery of urinary continence post-surgery.
\end{itemize}

\subsection*{Absolute Contraindications}
\begin{itemize}
  \item Distant metastatic disease: Evidence of cancer spread to distant organs, as radical prostatectomy is not curative in this setting.
  
  \item Severe comorbidities: Uncontrolled severe cardiac, pulmonary, or neurological conditions that render the patient an unacceptable surgical risk.
  
  \item Extremely short life expectancy: If the patient's life expectancy is less than 5 years due to other medical conditions, the risks may outweigh the benefits of surgery.
\end{itemize}

\subsection*{Relative Contraindications \& Precautions}
\begin{itemize}
  \item Advanced age: Very elderly patients may have increased surgical risks and slower recovery.
  
  \item Prior pelvic surgery or radiation: Previous surgeries or radiation can lead to significant scarring and increased complication risks.
  
  \item Large prostate gland: Extremely large prostates can complicate robotic or laparoscopic approaches, increasing operative time and blood loss.
  
  \item Active urinary tract infection: Surgery should be delayed until any active infection is treated to prevent complications.
  
  \item Severe coagulopathy: Uncorrectable bleeding disorders significantly increase surgical risk.
  
  \item Morbid obesity: Increases operative difficulty and risks of complications.
  
  \item Inflammatory bowel disease: May increase the risk of rectal injury during dissection.
\end{itemize}

\section*{Surgical Technique \& Procedure}
Radical prostatectomy is performed under general anesthesia. The patient is typically positioned in a supine or modified lithotomy position, depending on the surgical approach.

\begin{enumerate}
  \item **Anesthesia and Positioning:** General endotracheal anesthesia is administered. The patient is positioned supine, often with a steep Trendelenburg tilt for robotic or laparoscopic approaches to displace bowel superiorly.
  
  \item **Incision and Access:**
      \begin{itemize}
        \item **Open Radical Retropubic Prostatectomy (RRP):** A lower midline abdominal incision is made from the umbilicus to the pubic bone, allowing access to the retropubic space.
        \item **Robotic-Assisted Laparoscopic Prostatectomy (RALP):** 5-7 small incisions (ports) are made in the abdomen. Carbon dioxide gas is insufflated to create a pneumoperitoneum, providing a working space.
      \end{itemize}
  
  \item **Pelvic Lymph Node Dissection (PLND):** For higher-risk disease, a bilateral pelvic lymph node dissection is often performed first, removing lymph nodes for pathological staging.
  
  \item **Bladder Neck Dissection:** The bladder is separated from the prostate, and the bladder neck is incised circumferentially. The vas deferens and seminal vesicles are identified and ligated.
  
  \item **Neurovascular Bundle Preservation (Nerve-Sparing):** Meticulous dissection is performed along the posterolateral aspects of the prostate to identify and preserve the neurovascular bundles when oncologically safe.
  
  \item **Apex Dissection:** The prostate is separated from the external urethral sphincter and pelvic floor muscles at the apex to maximize continence preservation.
  
  \item **Prostate Removal:** The entire prostate gland and seminal vesicles are removed en bloc, and the specimen is sent for pathological examination.
  
  \item **Urethrovesical Anastomosis:** The remaining urethra is reconnected to the bladder neck using sutures, creating a watertight seal.
  
  \item **Drain Placement and Closure:** A surgical drain may be placed in the pelvis, and the incisions are closed in layers. A Foley catheter is left in the bladder to stent the anastomosis.
\end{enumerate}

\section*{Complications \& Risks}
Radical prostatectomy carries potential risks, categorized as general surgical complications and procedure-specific risks.

\begin{itemize}
  \item **General Surgical Complications:**
    \begin{itemize}
      \item **Bleeding:** Intraoperative blood loss requiring transfusion or postoperative hematoma formation.
      \item **Infection:** Surgical site infection or urinary tract infection.
      \item **Anesthesia Complications:** Adverse reactions to anesthetic agents, including respiratory depression and cardiac events.
      \item **Deep Vein Thrombosis (DVT) and Pulmonary Embolism (PE):** Blood clots that can lead to life-threatening complications.
      \item **Injury to adjacent organs:** Rarely, injury to the bowel or major blood vessels can occur.
    \end{itemize}
  
  \item **Procedure-Specific Risks:**
    \begin{itemize}
      \item **Urinary Incontinence:** Leakage of urine, which may improve over time with pelvic floor exercises.
      \item **Erectile Dysfunction:** Difficulty achieving or maintaining an erection, which may require medical assistance.
      \item **Bladder Neck Contracture:** Scar tissue formation leading to urinary obstruction.
      \item **Anastomotic Leak:** Leakage from the connection between the bladder and urethra.
      \item **Rectal Injury:** Accidental perforation of the rectum during dissection.
      \item **Lymphocele:** Collection of lymphatic fluid in the pelvis requiring drainage.
      \item **Nerve Injury:** Damage to pelvic nerves causing leg weakness or sensory changes.
      \item **Hernia:** Incisional hernia at the surgical site.
      \item **Shortening of the Penis:** Perceived or actual reduction in penile length post-surgery.
    \end{itemize}
\end{itemize}

\section*{Postoperative Care \& Recovery}
Following radical prostatectomy, patients are monitored in the Post-Anesthesia Care Unit (PACU) for recovery from anesthesia. Pain management is initiated, typically with intravenous analgesics, transitioning to oral medications as tolerated. A urinary catheter is placed to drain urine and allow healing of the urethrovesical anastomosis. Early mobilization is encouraged within the first 24-48 hours to prevent complications such as deep vein thrombosis.

Patients are usually discharged from the hospital within 1-3 days, depending on the surgical approach and individual recovery. Upon discharge, patients receive instructions for catheter care, pain management, and activity restrictions. The urinary catheter typically remains in place for 1-3 weeks, after which it is removed in an outpatient setting. Most men achieve social continence within 3-12 months, while recovery of erectile function may take 12-24 months and may require medical assistance.

During radical prostatectomy, the surgeon documents the appearance of the prostate, seminal vesicles, and surrounding tissues, noting any visible tumor extension or suspicious lymph nodes. The removed prostate gland and seminal vesicles are sent for pathological examination.

The pathology report provides critical information about the cancer, including:

\begin{itemize}
  \item **Gleason Score:** Indicates the aggressiveness of the cancer.
  \item **Tumor Volume and Location:** Describes the extent of cancer within the prostate.
  \item **Pathological Stage (pT stage):** Determines if the cancer was organ-confined or had extended beyond the prostate capsule.
  \item **Surgical Margins:** Presence of cancer cells at the edges of the removed tissue, indicating risk of recurrence.
  \item **Seminal Vesicle Invasion (SVI):** Indicates more advanced disease.
  \item **Lymph Node Status (pN stage):** Presence of cancer cells in removed lymph nodes.
  \item **Perineural Invasion:** Presence of cancer cells around nerves, indicating a higher risk of extracapsular extension.
\end{itemize}

These findings are crucial for accurate staging and predicting the patient's prognosis and need for further treatment.

A successful postoperative course following radical prostatectomy involves a predictable progression towards recovery. Initially, patients will experience incisional pain, managed effectively with analgesics. The presence of a Foley catheter for 1-3 weeks is standard. Upon catheter removal, most men will experience some degree of urinary leakage, which gradually improves with pelvic floor exercises. Social continence is typically achieved within 3-12 months.

From an oncologic perspective, patients expect an undetectable PSA level within 6-8 weeks after surgery, signifying successful cancer eradication. Recovery of erectile function is variable and may take 12-24 months, often requiring medical assistance. Patients can expect a gradual return to normal activities, with light activities resuming within a few weeks and full activity typically by 6-8 weeks post-surgery.

Postoperative care and long-term follow-up are essential for monitoring cancer recurrence and managing functional outcomes.

\begin{itemize}
  \item **Regular PSA Surveillance:** Serial PSA measurements are performed every 3-6 months for the first few years, then annually, to detect any rising PSA levels.
  
  \item **Postoperative Visits:** Scheduled appointments with the urologist are crucial to monitor recovery, discuss PSA results, and assess urinary and sexual function.
  
  \item **Catheter Removal:** The urinary catheter is typically removed 1-3 weeks post-surgery in an outpatient setting.
  
  \item **Pelvic Floor Physical Therapy:** Referral to a specialized physical therapist is often recommended to optimize recovery of urinary continence.
  
  \item **Management of Erectile Dysfunction:** Discussion and initiation of therapies such as oral medications or vacuum erection devices are common.
  
  \item **Activity Resumption:** Patients are guided on a gradual return to normal activities, with restrictions on heavy lifting for several weeks.
  
  \item **Warning Signs:** Patients are educated on symptoms requiring immediate medical attention, such as fever, severe pain, or inability to urinate.
\end{itemize}

Long-term follow-up continues for many years to monitor for late recurrences and manage any persistent side effects, ensuring ongoing quality of life.

\section*{Outcomes \& Prognosis}
Prostate cancer is the most common non-skin cancer among men in Western countries, with radical prostatectomy being one of the most frequently performed major surgeries for its treatment. The incidence of prostate cancer increases with age, with the median age at diagnosis around 66 years. Radical prostatectomy is most commonly performed in men aged 50-75 years who have clinically localized disease and a life expectancy of at least 10 years. 

Globally, prostate cancer incidence rates vary, with higher rates observed in North America and Western Europe. The choice of radical prostatectomy as a primary treatment varies geographically and depends on factors such as disease risk stratification and patient preference. Mortality directly attributable to radical prostatectomy is very low, typically less than 0.1-0.3\% in experienced centers. Morbidity, particularly related to urinary incontinence and erectile dysfunction, is more common, with rates varying significantly based on surgical technique and patient characteristics.

The outcomes of radical prostatectomy are generally excellent for appropriately selected patients with clinically localized prostate cancer, offering a high probability of long-term cancer control. For men with organ-confined disease, 10-year cancer-specific survival rates typically exceed 90-95\%. Even for those with more advanced localized disease, 10-year cancer-specific survival rates remain high, often in the range of 80-90\%. 

Biochemical recurrence rates vary depending on initial disease characteristics; for low- and intermediate-risk disease, 10-year biochemical recurrence-free survival can be 70-85\%. Prognostic factors significantly influencing outcomes include preoperative PSA level, biopsy Gleason score, and pathological findings from the resected specimen. While oncologic control is the primary goal, functional outcomes significantly impact quality of life. Most men experience some degree of urinary incontinence initially, with 60-90\% achieving social continence within a year. Recovery of erectile function is more challenging, with rates of potency recovery ranging from 20-80\%, heavily dependent on age and preoperative function. Landmark trials have provided valuable insights into the effectiveness of radical prostatectomy, demonstrating a survival benefit for surgery in certain patient cohorts.

\section*{References}
\begin{enumerate}
  \item MedlinePlus. Radical prostatectomy. U.S. National Library of Medicine; 2024. Available from: \url{https://medlineplus.gov/ency/article/007267.htm}.
  
  \item Wein AJ, Kavoussi LR, Partin AW, Peters CA, editors. Campbell-Walsh-Wein Urology. 12th ed. Philadelphia, PA: Elsevier; 2021.
  
  \item National Comprehensive Cancer Network (NCCN). NCCN Clinical Practice Guidelines in Oncology: Prostate Cancer. Version 1.2024. Available from: \url{https://www.nccn.org/professionals/physician_gls/pdf/prostate.pdf}.
  
  \item Walsh PC, Partin AW. Anatomic radical prostatectomy: evolution of the surgical technique. J Urol. 2004;172(5 Pt 2):S41-S46.
\end{enumerate}

\clearpage
